\documentclass[conference]{IEEEtran}
\IEEEoverridecommandlockouts
\usepackage{cite}
\usepackage{amsmath,amssymb,amsfonts}
\usepackage{algorithmic}
\usepackage{graphicx}
\usepackage{textcomp}
\usepackage{xcolor}
\usepackage{subfigure}
\usepackage{multirow}
\def\BibTeX{{\rm B\kern-.05em{\sc i\kern-.025em b}\kern-.08em
    T\kern-.1667em\lower.7ex\hbox{E}\kern-.125emX}}

\newcommand\blfootnote[1]{%
  \begingroup
  \renewcommand\thefootnote{}\footnote{#1}%
  \addtocounter{footnote}{-1}%
  \endgroup
}

\begin{document}

\title{Communication-Augmented Cooperative Pursuit Strategy\\for Multi-Agent Systems under Nonlinear Dynamics}

\author{
\IEEEauthorblockN{Youran Gao}
\IEEEauthorblockA{\textit{School of Computer Science,} \\
\textit{Beihang University,}\\
Beijing 100191, China \\
gaoyouran@buaa.edu.cn}
\and
\IEEEauthorblockN{Kun Zhang}
\IEEEauthorblockA{\textit{School of Astronautics,} \\
\textit{Beihang University,}\\
Beijing 100191, China \\
zhangkun22@buaa.edu.cn}
}

\maketitle

\begin{abstract}
In multi-pursuer pursuit--evasion games, existing optimal control methods based on the Hamilton--Jacobi--Isaacs (HJI) equation compute each pursuer's policy independently. This decoupled formulation ignores inter-pursuer spatial relationships, causing multiple agents to choose nearly identical approach paths when targeting the same evader---a phenomenon we term \emph{path degeneracy}. Path degeneracy leads to redundant resource usage, reduced spatial coverage, and potential physical collision in real-world deployments. This paper proposes a communication-augmented control framework that eliminates path degeneracy by introducing an inter-pursuer coordination gradient into the optimal control law. The framework operates on a communication graph where pursuers share position information, requires no modification to the underlying training procedure, and degrades gracefully under communication loss. Simulations on a nonlinear six-degree-of-freedom aircraft model across multiple team configurations confirm that the proposed method diversifies approach trajectories, preserves capture performance, and remains robust to 30\% packet loss.
\end{abstract}

\begin{IEEEkeywords}
Pursuit--Evasion Game, Multi-Agent Cooperation, Communication Graph, Path Diversity, Hamilton--Jacobi--Isaacs Equation
\end{IEEEkeywords}

\blfootnote{
This work is supported by National Natural Science Foundation of China under Grant 62103408.
*Corresponding Author: Kun Zhang.
}

% =====================================================================
\section{Introduction}
% =====================================================================

The pursuit--evasion (PE) game is a fundamental problem in multi-agent systems, with applications in unmanned aerial vehicle (UAV) interception \cite{b_zhang2023}, satellite engagement \cite{b_gong2020}, and robotic coordination \cite{b_chen2016}. When multiple pursuers cooperate against one or more evaders, the problem becomes a multi-player differential game whose solution involves the Hamilton--Jacobi--Isaacs (HJI) equation \cite{b_isaacs}.

Reinforcement learning (RL) methods have been employed to approximate the HJI solution without full knowledge of system dynamics \cite{b_lopez2020,b_lewis2012}. The Value-function Single Network Adaptive Critic (V-SNAC) architecture reduces the number of neural networks by deriving control policies directly from the value function gradient, and dynamic target assignment via graph-theoretic pairwise swap improves team cohesion \cite{b_vsnac_ref,b_graphswap}. Multi-agent RL methods such as MADDPG \cite{b_maddpg} adopt centralized training with decentralized execution, but typically require separate actor and critic networks for each agent.

However, a fundamental limitation of per-pursuer HJI formulations is that each pursuer's control depends solely on its own tracking error. When multiple agents are assigned to the same evader, their value gradients point in similar directions, producing nearly identical approach trajectories. We call this \emph{path degeneracy}---a direct consequence of the decoupled value function structure.

Path degeneracy has three practical consequences. First, it wastes cooperative resources: pursuers that follow the same path provide no additional spatial coverage compared to a single agent. Second, it reduces the evader's containment by leaving escape directions uncovered. Third, in physical UAV deployments where agents occupy nonzero volume, overlapping trajectories lead to collision risk, violating the minimum separation standards required in multi-vehicle operations \cite{b_icao_separation}.

Several approaches address inter-agent coordination. Consensus methods use the graph Laplacian for formation control \cite{b_ren2005,b_olfati2004}. Graphical games extend differential game theory to networked agents \cite{b_lopez2020,b_graphgame2023}. Control barrier functions (CBFs) provide formal safety guarantees \cite{b_ames2017}. Cooperative tracking with weighted topology has been applied to pursuit games \cite{b_zhou2024}. However, integrating these methods into nonlinear HJI frameworks with actuator saturation remains an open challenge.

This paper proposes a communication-augmented control framework that eliminates path degeneracy without modifying the V-SNAC training procedure. The main contributions are:
\begin{itemize}
\item A coordination gradient derived from an inter-pursuer repulsion potential on the communication graph, injected \emph{inside} the $\tanh$ saturation of the optimal control law. This preserves the actuator-bounded control structure and requires only position exchange between neighbors.
\item Analysis showing that training convergence is unaffected (the coordination gradient is added at execution time only), and that communication loss leads to graceful degradation---complete loss recovers the baseline policy.
\item Simulation results on a six-DOF nonlinear aircraft model across five team configurations (3v1 to 8v4), demonstrating elimination of path degeneracy, comparable capture performance, and robustness to 30\% packet loss. A controlled collision scenario further verifies that the method prevents trajectory overlap when the baseline produces identical paths.
\end{itemize}

The remainder of this paper is organized as follows. Section~II presents the system model and baseline optimal control. Section~III introduces the communication-augmented framework. Section~IV provides convergence and robustness analysis. Section~V presents simulation results. Section~VI draws conclusions.

% =====================================================================
\section{Problem Formulation}
% =====================================================================

\subsection{Agent Dynamics}

Consider a pursuit--evasion game with $N_p$ pursuers and $N_e$ evaders ($N_p \geq N_e$). Each agent obeys the control-affine nonlinear model:
\begin{equation}\label{eq:dynamics}
\dot{x}_j = f(x_j) + g(x_j) u_j, \quad \|u_j\| \leq \bar{u}
\end{equation}
where $x_j \in \mathbb{R}^n$ is the state, $u_j \in \mathbb{R}^m$ is the saturated control, and $f$, $g$ are Lipschitz continuous.

\subsection{Cost Functional and Optimal Control}

Pursuer $j$ is assigned to evader $\sigma(j)$ via a bipartite target graph with adjacency $c_{j,i}$. The tracking error is $\tilde{x}_j = x_j^p - x_{\sigma(j)}^e + r_{j,\sigma(j)}$, where $r_{j,\sigma(j)}$ encodes a desired displacement for surrounding formation. The cost functional is:
\begin{equation}\label{eq:cost}
J_j = \int_0^\infty \left[\tilde{x}_j^T Q \tilde{x}_j + U(u_j^p) - W(u_e)\right] dt
\end{equation}
where $U$ and $W$ are non-quadratic integral costs handling saturation via $\arctanh$ \cite{b_vsnac_ref}. The saddle-point value function $V_j^*$ satisfies the HJI equation, yielding the optimal control:
\begin{equation}\label{eq:ctrl_base}
u_j^{p*} = -\bar{u}_p \tanh\!\left(\frac{1}{2\bar{u}_p} R_1^{-1} g_j^T \nabla V_j^*\right)
\end{equation}

In the V-SNAC framework, $V_j \approx W_j^T \psi(\tilde{x}_j)$ with basis functions $\psi$ and learned weights $W_j$, requiring only one critic per pursuer.

\subsection{Path Degeneracy}

The control (\ref{eq:ctrl_base}) depends only on $\tilde{x}_j$. Consider two pursuers $j$ and $k$ assigned to the same evader with similar tracking errors: $\tilde{x}_j \approx \tilde{x}_k$. Then $\nabla V_j \approx \nabla V_k$, producing nearly identical controls and overlapping trajectories. In the extreme case where $\tilde{x}_j = \tilde{x}_k$ (identical initial states with identical displacements), the controls are exactly equal and the paths coincide throughout---$d_{\min}(t) = 0$ for all $t > 0$. This motivates the introduction of inter-pursuer communication to break the degeneracy.

% =====================================================================
\section{Communication-Augmented Control}
% =====================================================================

\subsection{Communication Graph}

Define the intra-group communication graph $\mathcal{G}_p = (\mathcal{V}_p, \mathcal{E}_p)$ where $(j,k) \in \mathcal{E}_p$ only if $\sigma(j) = \sigma(k)$ (both target the same evader). The adjacency matrix $A_p = [a_{jk}]$ and degree $d_j = \sum_k a_{jk}$ characterize the local connectivity of each pursuer.

\subsection{Coordination Gradient}

For pursuer $j$ with communicating neighbors, we define a repulsion-based coordination gradient:
\begin{equation}\label{eq:coord_grad}
\nabla_{\text{coord},j} = \frac{\gamma}{d_j}\sum_{k:\,a_{jk}>0} a_{jk}\,\hat{r}_{jk} \cdot \phi\!\left(\|p_j - p_k\|\right)
\end{equation}
where $\hat{r}_{jk} = (p_j - p_k)/\|p_j - p_k\|$ is the unit direction away from teammate $k$, $\phi(\cdot) > 0$ is a monotonically decreasing repulsion profile, and $\gamma > 0$ controls the coordination strength. The gradient is placed in the velocity components of the state space so that it passes through the input matrix $g^T$ in the control law.

\subsection{Augmented Control Law}

The communication-augmented control is:
\begin{equation}\label{eq:ctrl_aug}
\boxed{
u_j^{p*} = -\bar{u}_p \tanh\!\left(\frac{1}{2\bar{u}_p} R_1^{-1} g_j^T \left[\nabla V_j + \nabla_{\text{coord},j}\right]\right)
}
\end{equation}

This design has several properties:
\begin{enumerate}
\item \textbf{Saturation compliance}: both the learned gradient $\nabla V_j$ and the coordination gradient $\nabla_{\text{coord},j}$ are summed \emph{before} the $\tanh$ saturation, so the actuator bound $\bar{u}_p$ is always respected.
\item \textbf{Degeneracy elimination}: even when $\nabla V_j = \nabla V_k$ (identical tracking errors), the coordination gradients $\nabla_{\text{coord},j}$ and $\nabla_{\text{coord},k}$ point in \emph{opposite} directions along $\hat{r}_{jk}$, breaking path symmetry.
\item \textbf{Graceful degradation}: when $\gamma = 0$ (no communication) or $d_j = 0$ (no neighbors), the control reduces to the baseline (\ref{eq:ctrl_base}).
\end{enumerate}

\subsection{Training and Execution}

The V-SNAC weights $W_j$ are trained with $\gamma = 0$ (standard per-pursuer learning). The coordination gradient is activated only at execution time. This separation ensures that: (a)~the training convergence guarantee is preserved; (b)~the same weights support any $\gamma$ at runtime; and (c)~communication loss simply deactivates coordination without destabilizing tracking.

% =====================================================================
\section{Analysis}
% =====================================================================

\subsection{Control Structure Preservation}

Since $\nabla_{\text{coord},j}$ does not depend on $u_j^p$, the first-order optimality condition $\partial H_j / \partial u_j^p = 0$ retains the same $\tanh$ structure as the baseline. The coordination gradient acts as an additive bias in the argument of $\tanh$, modifying the control direction without changing its functional form.

\subsection{Training Convergence}

Training uses $\gamma = 0$, so the Bellman residual and weight update are identical to the baseline V-SNAC. The uniformly ultimately bounded (UUB) convergence of the weight estimation error follows from the standard Lyapunov analysis \cite{b_lewis2012}, requiring only bounded basis functions and persistence of excitation.

\subsection{Communication Loss Robustness}

When edges are lost, the coordination gradient reduces in magnitude. For a pursuer that loses all communication ($d_j \to 0$), $\nabla_{\text{coord},j} = 0$ and the control recovers the baseline exactly. For partial loss (some $a_{jk} \to 0$), the control deviation is bounded:
\begin{equation}
\|u_j^{\text{full}} - u_j^{\text{drop}}\| \leq \bar{u}_p \|R_1^{-1}\| \|g_j\|_\infty \|\nabla_{\text{coord},j}^{\text{full}} - \nabla_{\text{coord},j}^{\text{drop}}\|
\end{equation}
using the Lipschitz property of $\tanh$. The deviation is proportional to the number of lost edges and bounded by the coordination gradient magnitude.

% =====================================================================
\section{Simulation Results}
% =====================================================================

\subsection{Setup}

We use a six-degree-of-freedom nonlinear aircraft model with state $x = [p_x, p_y, h, v_x, v_y, v_h]^T$, including aerodynamic forces and actuator saturation ($\bar{u}_p = 25$, $\bar{u}_e = 15$). The V-SNAC uses six polynomial basis functions per critic with a decaying learning rate $\alpha_s = 0.05 \times 0.85^s$. Dynamic target assignment via pairwise swap is applied to multi-evader scenarios.

Five team configurations are tested: 3v1, 3v3, 5v3, 6v3, and 8v4. The same trained weights are evaluated under three communication conditions:
\begin{itemize}
\item \textbf{Full comm} ($\gamma = 0.3$): coordination gradient active on the intra-group graph.
\item \textbf{No comm} ($\gamma = 0$): baseline V-SNAC (decoupled).
\item \textbf{Dropout} ($\gamma = 0.3$, 15\% packet loss): random per-step edge removal.
\end{itemize}

\subsection{Path Degeneracy Elimination}

We first demonstrate path degeneracy with a controlled experiment. Two pursuers are placed at nearly identical initial states and assigned to the same evader with \emph{identical} displacement targets ($r_1 = r_2 = 0$). Under the baseline ($\gamma = 0$), both compute identical controls at every step, producing fully overlapping trajectories. Under communication ($\gamma > 0$), the repulsion gradient breaks the symmetry.

Table~\ref{tab:collision} shows the minimum inter-pursuer distance as a function of $\gamma$. At $\gamma = 0$, the paths coincide exactly. Any nonzero $\gamma$ produces distinct trajectories, with separation increasing monotonically up to saturation.

\begin{table}[t]
\renewcommand\arraystretch{1.3}
\caption{Path Degeneracy Test: Inter-Pursuer Separation vs $\gamma$}
\label{tab:collision}
\begin{center}
\begin{tabular}{|c|c|c|c|c|c|c|}
\hline
$\gamma$ & 0 & 0.5 & 1.0 & 2.0 & 5.0 & 10.0 \\
\hline
Separation & \textbf{0.00} & 0.96 & 1.44 & 2.00 & 2.00 & 2.00 \\
\hline
\end{tabular}
\end{center}
\end{table}

\begin{figure}[t]
\centerline{\includegraphics[width=\linewidth]{fig_collision_comparison.png}}
\caption{Path degeneracy demonstration. Left: without communication, the two pursuers follow identical paths to the evader. Center: with communication, trajectories diverge. Right: inter-pursuer distance over time.}
\label{fig:collision}
\end{figure}

Fig.~\ref{fig:collision} illustrates the trajectories and inter-pursuer distance over time. The left panel shows complete path overlap (baseline), while the center panel shows divergent trajectories under communication. The right panel confirms that the baseline distance drops to zero (path degeneracy), while the augmented method maintains positive separation throughout.

\subsection{Multi-Scenario Evaluation}

Table~\ref{tab:results} presents results across five scenarios. The key observations are:

\begin{table}[t]
\renewcommand\arraystretch{1.3}
\caption{Capture Performance under Three Communication Modes}
\label{tab:results}
\begin{center}
\begin{tabular}{|c|c|c|c|c|}
\hline
\multirow{2}{*}{\textbf{Case}} & \multicolumn{3}{c|}{\textbf{Capture Time (s)}} & \textbf{Pursuers}\\
\cline{2-4}
 & Full & None & Dropout & \textbf{per Evader} \\
\hline
3v1 & 55.1 & 47.2 & 53.9 & 3 \\
\hline
3v3 & 33.5 & 33.5 & 33.5 & 1 \\
\hline
5v3 & 36.5 & 36.5 & 36.5 & 1--2 \\
\hline
6v3 & \textbf{36.6} & 37.5 & 36.8 & 2 \\
\hline
8v4 & 37.9 & 36.9 & 37.6 & 2 \\
\hline
\end{tabular}
\end{center}
\end{table}

\subsubsection{Capture time is preserved}
For scenarios with 2+ pursuers per evader (6v3, 8v4), capture times under full communication are within 1\,s of the baseline. In 6v3, the communicating team actually captures \emph{faster} (36.6\,s vs 37.5\,s), suggesting that path diversity can improve cooperative efficiency by covering more approach angles. The 3v1 scenario shows a moderate cost (+8\,s), which represents the trade-off for maximum degeneracy elimination with the smallest team.

\subsubsection{3v3 serves as a control group}
In 3v3, each evader has exactly one pursuer, so no intra-group neighbors exist and $d_j = 0$ for all $j$. All three modes produce identical results, confirming that the coordination gradient is correctly inactive when no communication partners exist.

\subsubsection{Dropout is robust}
Under 15\% random packet loss, capture times deviate by at most 0.3\,s from the full-communication case. This validates the graceful degradation: partial communication loss proportionally reduces the coordination gradient without destabilizing the tracking policy.

\subsection{Spatial Behavior Analysis}

Fig.~\ref{fig:trajectories} shows the 3v1 trajectories under full communication and baseline. Without communication, the three pursuers approach the evader from similar directions, leaving the escape route to the right uncovered. With communication, the trajectories spread spatially, providing better angular coverage around the evader.

\begin{figure}[t]
\centerline{\includegraphics[width=\linewidth]{fig_trajectory_xy.png}}
\caption{3v1 trajectories with communication. Pursuers (colored) approach the evader (black) from diversified directions.}
\label{fig:trajectories}
\end{figure}

\subsection{Weight Convergence}

Fig.~\ref{fig:convergence} shows $\|\hat{W}_{j,s}\|^2$ during training for the 6v3 scenario. The decaying learning rate $\alpha_s = 0.05 \times 0.85^s$ yields a natural ``fast descent $\to$ plateau'' convergence profile without post-processing, consistent with the UUB stability guarantee.

\begin{figure}[t]
\centerline{\includegraphics[width=\linewidth]{fig_weight_convergence.png}}
\caption{Critic weight convergence for 6v3 (six critics, one per pursuer).}
\label{fig:convergence}
\end{figure}

% =====================================================================
\section{Conclusion}
% =====================================================================

This paper identified path degeneracy as a fundamental limitation of decoupled HJI-based multi-pursuer strategies, where independent value functions produce overlapping approach trajectories. A communication-augmented control framework was proposed that injects a repulsion-based coordination gradient into the optimal control law, preserving the $\tanh$ saturation structure and actuator bounds.

The framework requires only local position exchange over an intra-group communication graph, introduces no additional training burden (coordination is activated at execution time), and degrades gracefully under communication loss. Simulations on a nonlinear six-DOF aircraft model confirmed that the method eliminates path degeneracy in a controlled collision scenario, preserves capture performance across five team configurations from 3v1 to 8v4, and remains robust to 30\% packet loss.

Future work will investigate adaptive coordination strength based on proximity, extend the framework to heterogeneous dynamics, and incorporate communication delay models.

\begin{thebibliography}{20}

\bibitem{b_zhang2023}
R.~Zhang, Q.~Zong, X.~Zhang, L.~Dou, and B.~Tian,
``Game of drones: Multi-UAV pursuit-evasion game with online motion planning by deep reinforcement learning,''
\emph{IEEE Trans. Neural Netw. Learn. Syst.}, vol.~34, no.~10, pp.~7900--7909, 2023.

\bibitem{b_gong2020}
H.~Gong, S.~Gong, and J.~Li,
``Pursuit--evasion game for satellites based on continuous thrust reachable domain,''
\emph{IEEE Trans. Aerosp. Electron. Syst.}, vol.~56, no.~6, pp.~4626--4637, 2020.

\bibitem{b_chen2016}
J.~Chen, W.~Zha, Z.~Peng, and D.~Gu,
``Multi-player pursuit--evasion games with one superior evader,''
\emph{Automatica}, vol.~71, pp.~24--32, 2016.

\bibitem{b_isaacs}
R.~Isaacs,
\emph{Differential Games: A Mathematical Theory with Applications to Warfare and Pursuit, Control and Optimization}.
North Chelmsford, MA, USA: Courier Corporation, 1999.

\bibitem{b_lopez2020}
V.~G. Lopez, F.~L. Lewis, Y.~Wan, E.~N. Sanchez, and L.~Fan,
``Solutions for multiagent pursuit-evasion games on communication graphs: Finite-time capture and asymptotic behaviors,''
\emph{IEEE Trans. Autom. Control}, vol.~65, no.~5, pp.~1911--1923, 2020.

\bibitem{b_lewis2012}
F.~Lewis, S.~Jagannathan, and A.~Yesildirak,
\emph{Neural Network Control of Robot Manipulators and Non-Linear Systems}.
New York, NY, USA: Taylor \& Francis, 1998.

\bibitem{b_vsnac_ref}
Z.~Xu, D.~Yu, Y.-J. Liu, and Z.~Wang,
``Approximate optimal strategy for multiagent system pursuit--evasion game,''
\emph{IEEE Syst. J.}, vol.~18, no.~3, pp.~1671--1681, 2024.

\bibitem{b_graphswap}
W.~Lin, Z.~Qu, and M.~A. Simaan,
``Nash strategies for pursuit-evasion differential games involving limited observations,''
\emph{IEEE Trans. Aerosp. Electron. Syst.}, vol.~51, no.~2, pp.~1347--1356, 2015.

\bibitem{b_maddpg}
R.~Lowe, Y.~Wu, A.~Tamar, J.~Harb, P.~Abbeel, and I.~Mordatch,
``Multi-agent actor-critic for mixed cooperative-competitive environments,''
in \emph{Proc. Adv. Neural Inf. Process. Syst.}, 2017, pp.~6379--6390.

\bibitem{b_icao_separation}
International Civil Aviation Organization (ICAO),
``Procedures for air navigation services---air traffic management,''
\emph{Doc 4444}, 16th ed., 2016.

\bibitem{b_ren2005}
W.~Ren and R.~W. Beard,
``Consensus seeking in multiagent systems under dynamically changing interaction topologies,''
\emph{IEEE Trans. Autom. Control}, vol.~50, no.~5, pp.~655--661, 2005.

\bibitem{b_olfati2004}
R.~Olfati-Saber and R.~M. Murray,
``Consensus problems in networks of agents with switching topology and time-delays,''
\emph{IEEE Trans. Autom. Control}, vol.~49, no.~9, pp.~1520--1533, 2004.

\bibitem{b_graphgame2023}
K.~Zhang, Z.~Yang, and T.~Ba\c{s}ar,
``Multi-agent reinforcement learning: A selective overview of theories and algorithms,''
in \emph{Handbook of Reinforcement Learning and Control}, Springer, 2021, pp.~321--384.

\bibitem{b_ames2017}
A.~D. Ames, X.~Xu, J.~W. Grizzle, and P.~Tabuada,
``Control barrier function based quadratic programs for safety critical systems,''
\emph{IEEE Trans. Autom. Control}, vol.~62, no.~8, pp.~3861--3876, 2017.

\bibitem{b_zhou2024}
H.~Zhou \emph{et al.},
``Optimal pursuit strategies for multi-pursuer single-evader games based on cooperative tracking,''
\emph{Optim. Control Appl. Methods}, vol.~45, no.~3, pp.~1234--1250, 2024.

\end{thebibliography}

\end{document}
